
\section{Large-data exact unforced controls and the nonaxisymmetric angular-momentum escape gate}
\label{sec:xxxv-nonaxis-torque}

Edition~XXII already established that the bounded-circulation
maximum principle prevents an axisymmetric unforced solution from
retaining the particular growing azimuthal core of the original
forced source \cite{openai2026}. We now test the distinct
large-initial-data proposal with an exact no-swirl solution
and derive an additional sharp quantitative requirement on
\emph{nonaxisymmetric} transport. Neither result solves the
unforced three-dimensional Navier--Stokes regularity problem.

\begin{theorem}[A large-critical-norm nonzero datum with an exact globally regular unforced flow]
\label{thm:xxxv-large-data}
For any $\nu>0$ set $r=\sqrt{x^2+y^2}$ and
$\rho^2=r^2+z^2$. Choose the explicit whole-space datum
\begin{equation}\label{eq:xxxv-initial}
\boxed{\quad
v_0(x,y,z)=2\nu e^{-\rho^2}
\bigl(xz,yz,1-r^2\bigr).\quad}
\end{equation}
It is nonzero, smooth, Schwartz, exactly divergence-free
and axisymmetric \emph{without swirl}. In the unitary
Fourier convention and the critical space
$\|a\|_{\mathcal X^{-1}}=\int|\xi|^{-1}|\widehat a(\xi)|\,d\xi$,
\begin{equation}\label{eq:xxxv-large-critical}
\boxed{\quad
\|v_0\|_{\mathcal X^{-1}}=\sqrt2\,\pi^2\nu,\qquad
\frac{(2\pi)^{-3/2}}{\nu}\|v_0\|_{\mathcal X^{-1}}
=\frac{\sqrt\pi}{2}>\frac1{16}.\quad}
\end{equation}
The classical global axisymmetric no-swirl theorem
(e.g. Gallay--\v{S}ver\'ak \cite{gallaySverak2015})
gives a unique global smooth exactly \emph{unforced}
finite-energy solution $v$ arising from this datum,
with $v_\theta\equiv0$. Its initial azimuthal vorticity is
\[
\omega_\theta(r,z,0)=
2\nu r(5-2\rho^2)e^{-\rho^2},\qquad
\left\|\frac{\omega_\theta(0)}r\right\|_\infty=10\nu.
\]
For comparison with the published viscosity-one forced source
$u_B$ of \cite{openai2026}, take $\nu=1$; at other
viscosities one must first use a correspondingly normalized
forced background. In the viscosity-one case, the correction $w=v-u_B$ satisfies the complete
zero-force correction equation and
\begin{equation}\label{eq:xxxv-swirl-cancel}
\boxed{\qquad w_\theta=-u_{B,\theta}\quad
\text{for every }0\le t<1.\qquad}
\end{equation}
Consequently this \emph{large-data} exact unforced
comparison cancels the forced azimuthal singularity
identically.
\end{theorem}
\begin{proof}
Let $\psi=e^{-\rho^2}$ and $e_z=(0,0,1)$.
Then $v_0=(\nu/2)\operatorname{curl}
\operatorname{curl}(\psi e_z)$, giving
$\nabla\cdot v_0=0$ and the displayed
axisymmetric no-swirl components
$v_r=2\nu rz e^{-\rho^2}$,
$v_z=2\nu(1-r^2)e^{-\rho^2}$.
In Fourier space,
\[
\widehat v_0(\xi)=
\frac{\nu}{2^{5/2}}e^{-|\xi|^2/4}
\bigl(|\xi|^2e_z-\xi_z\xi\bigr).
\]
Its modulus is
$\nu 2^{-5/2}|\xi|^2\sin\varphi
e^{-|\xi|^2/4}$, where $\varphi$ is the angle
with the $z$ axis. The exact angular and radial
integrals are $\int_{\mathbb S^2}\sin\varphi\,d\Omega=\pi^2$
and $\int_0^\infty q^3e^{-q^2/4}dq=8$.
This proves \eqref{eq:xxxv-large-critical}.
Direct cylindrical differentiation gives the vorticity formula;
the ratio obeys the no-swirl advection--diffusion
equation with the lifted meridional Laplacian, and its
initial supremum is $10\nu$.
Classical no-swirl global regularity supplies $v(t)$ and
preserves $v_\theta=0$ for all time.
The difference $w=v-u_B$ obeys
\[
\partial_tw-\nu\Delta w+
\mathbb P\nabla\cdot
(u_B\otimes w+w\otimes u_B+w\otimes w)
=-\mathbb P f_B,
\]
where $f_B$ is the \emph{complete} Section~10 physical force.
The full unforced equation of $v$ holds with its reconstructed
pressure; \eqref{eq:xxxv-swirl-cancel} follows identically.
\end{proof}

\begin{theorem}[Exact nonaxisymmetric angular-momentum transport identity]
\label{thm:xxxv-angular-mean}
Let $u$ be any smooth divergence-free solution of the
\emph{unforced} Navier--Stokes equations on $\mathbb R^3$
up to a time $T>0$. Let
$\Gamma=r u_\theta$ and denote averages over
$\theta\in[0,2\pi]$ by a bar and fluctuations by
primes, $f'=f-\bar f$. Then, wherever $r>0$,
\begin{equation}\label{eq:xxxv-averaged-gamma}
\boxed{
(\partial_t+\bar u_r\partial_r+\bar u_z\partial_z)
\bar\Gamma=
\nu(\partial_{rr}-r^{-1}\partial_r+\partial_{zz})
\bar\Gamma+\mathcal T,
}
\end{equation}
where the \emph{entire} nonaxisymmetric transport term is
\begin{equation}\label{eq:xxxv-torque}
\boxed{\qquad
\mathcal T=
-\frac1r\partial_r(r\,\overline{u'_r\Gamma'})
-\partial_z\overline{u'_z\Gamma'}.\qquad}
\end{equation}
There is no discarded mean pressure torque:
$\overline{\partial_\theta p}=0$ identically.
\end{theorem}
\begin{proof}
For the Cartesian angular momentum
$\Gamma=xu_y-yu_x$, the cylindrical momentum equations
give, with the full material derivative $D_t$,
\[
D_t\Gamma=
\nu(\Delta-2r^{-1}\partial_r)\Gamma+
2\nu r^{-1}\partial_\theta u_r-\partial_\theta p.
\]
Because $\nabla\cdot u=0$, the advection has the
conservative form
\[
u\cdot\nabla\Gamma
=r^{-1}\partial_r(ru_r\Gamma)+
r^{-1}\partial_\theta(u_\theta\Gamma)+
\partial_z(u_z\Gamma).
\]
Average over $\theta$. The two explicit angular
derivatives and the angular Laplacian vanish, giving
\[
\partial_t\bar\Gamma+
r^{-1}\partial_r(r\,\overline{u_r\Gamma})+
\partial_z\overline{u_z\Gamma}
=\nu(\partial_{rr}-r^{-1}\partial_r+\partial_{zz})
\bar\Gamma.
\]
The mean meridional velocity satisfies
$r^{-1}\partial_r(r\bar u_r)+\partial_z\bar u_z=0$.
Decompose the mean products into
$\bar u_r\bar\Gamma+\overline{u'_r\Gamma'}$
and $\bar u_z\bar\Gamma+\overline{u'_z\Gamma'}$,
giving \eqref{eq:xxxv-averaged-gamma}.
\end{proof}

\begin{theorem}[Quantitative nonaxisymmetric escape dichotomy]
\label{thm:xxxv-dichotomy}
Assume in addition that $\bar\Gamma$ has bounded initial
data and satisfies the standard maximum-principle hypotheses
at the symmetry axis and spatial infinity, with
$M_0=\|\bar\Gamma(0)\|_\infty<\infty$.
Suppose the completed unforced velocity
retains at least one-quarter of the \emph{positive}
source azimuthal amplitude at the trajectory
$x_\tau=(r_\tau,0,z_\tau)$ of (10.20)--(10.21):
\[
r_\tau=\sqrt{2X_{\rm in}\tau},\quad
u_\theta(r_\tau,0,z_\tau,1-\tau)
\ge(e_0/4)\tau^{-1/2-h},
\quad e_0>0,\ h>0.
\]
Then, for all sufficiently small $\tau>0$,
\begin{equation}\label{eq:xxxv-escape}
\boxed{\begin{aligned}
&\int_0^{1-\tau}
 \|\mathcal T(s)\|_{L^\infty_{r,z}}\,ds\\
&\qquad+\sqrt{\frac{\pi}{6}}\,
 \|\partial_\theta\Gamma(r_\tau,\cdot,z_\tau,1-\tau)\|_{L^2_\theta}
\ge \frac{\sqrt{2X_{\rm in}}e_0}{4}\tau^{-h}-M_0.
\end{aligned}}
\end{equation}
In particular, if \emph{both} quantities on the left
are $o(\tau^{-h})$, such source-swirl retention is impossible.
The coefficient $\sqrt{\pi/6}$ is the optimal
point-evaluation constant for the periodic zero-mean
homogeneous $H^1$ angular inequality, not a PDE
sharpness claim for \eqref{eq:xxxv-escape}.
\end{theorem}
\begin{proof}
The inhomogeneous scalar maximum principle applied to
\eqref{eq:xxxv-averaged-gamma} gives
\[
\|\bar\Gamma(t)\|_\infty\le
M_0+\int_0^t\|\mathcal T(s)\|_\infty ds.
\]
For a periodic $H^1$ function $g(\theta)$, Fourier series
and Cauchy--Schwarz yield
\[
|g(0)-\bar g|
\le \left(
 \frac1{2\pi}\sum_{n\ne0}\frac1{n^2}
 \right)^{1/2}\|g'\|_{L^2_\theta}
=\sqrt{\frac{\pi}{6}}\|g'\|_{L^2_\theta}.
\]
Indeed $\sum_{n\ne0}n^{-2}=\pi^2/3$;
truncations of the Fourier sequence $g_n=n^{-2}$
approach equality. Apply this inequality to
$g(\theta)=\Gamma(r_\tau,\theta,z_\tau,1-\tau)$,
then use the retained amplitude and
$r_\tau\tau^{-1/2-h}
=\sqrt{2X_{\rm in}}\tau^{-h}$.
\end{proof}

\paragraph{An explicit nonzero Reynolds angular-momentum flux.}
On an annular plateau choose the $z$ component
of a smooth compactly supported vector potential
as $A_z=r\cos\theta+r^3\sin\theta$.
Set $u=\operatorname{curl}(A_ze_z)$, extending
with radial and axial smooth cutoffs that equal one
on the plateau. Then $u$ is exactly divergence-free.
Direct calculation at the plateau gives
\[
u_r=-\sin\theta+r^2\cos\theta,\qquad
u_\theta=-\cos\theta-3r^2\sin\theta,
\quad\bar\Gamma=0,
\]
\[
\overline{u'_r\Gamma'}=r^3,\qquad
-\frac1r\partial_r(r\overline{u'_r\Gamma'})=-4r^2.
\]
Thus nonaxisymmetric Reynolds transport can be
nonzero for a smooth solenoidal field. This is a
\emph{kinematic flux example}, not a solution of
the time-dependent unforced PDE or a construction
of the required divergent torque.

\paragraph{Proof boundary.}
The original Section~10 source provides the
positive shrinking-core swirl, while the
no-swirl global regularity theorem and periodic
angular trace inequality are classical prior art.
The new source-targeted condition is the exact
flux/oscillation dichotomy \eqref{eq:xxxv-escape};
it does not produce a nonaxisymmetric unforced
solution with those costs, and the source's
numerical profile still lacks a frozen certificate.
All six large-data Gaussian formulas and the
nonaxisymmetric flux identities were audited
independently with symbolic algebra and rational
fixtures. A separate Lean module proves only
scalar algebra; the full Navier--Stokes PDE,
maximum principle and existence arguments are
not kernel formalized.
